\chapter{Discussion}
\label{ch:discussion}

\todonote{Written after ch07. Planned structure below.}

\section{Synthesis across research questions}
% RQ1-RQ5 answered explicitly, each pointing at the figures/tables that carry
% the evidence; where results are null, say so plainly.

\section{Limitations and threats to validity}
% - Scale: <= 20 qubits, everything classically simulable (by design).
% - Non-nested tuning: residual overlap between the tuning fold's training
%   data and other folds' evaluations; direction and rough size of the bias.
% - Budget sensitivity: 50-trial ceiling may favor models with smaller search
%   spaces; mitigated by the "size philosophy" cap, not eliminated.
% - Dataset scope: standard benchmarks + engineered synthetics; no natural
%   quantum data.
% - Hardware sample: few QPU runs, one vendor, one calibration window.

\section{When are quantum or hybrid approaches worth considering?}
% Decision-oriented reading of the evidence: task structure (parity-like
% interactions), data regime (small n), cost tolerance, and the honest
% competitiveness picture at NISQ scale.
